\section{All levels of Hastings' hierarchical projection}
\label{sec:hastings05_hierarchy_levels}

% Hastings, arXiv:cond-mat/0508554v3, proj3new.tex, lines 668--707.

The hierarchy of~\cite[lines 668--707]{Hastings2006Solving} defines level-$m+1$
projections $P_k^{m+1}$ onto the spectral window of the pair operator $M_k^m$
inside $\operatorname{ran}(P_{2k}^mP_{2k+1}^m)$. Measuring each deficiency
$\lVert\Psi_a-P_k^m\Psi_a\rVert$ against $\Psi_a$ directly yields a recurrence
for squared deficiencies with prefactor $3^m-1$ in place of
$2^m-1$~\cite{gap:hastings05_hierarchical_projection}.

\begin{lemma}[One level against the ground state]
    \label{thm:hastings05_level_against_ground_state}
    \lean{TNLean.MPS.AreaLaw.LocalSolution.norm_sub_compressedWindowProjection_mul_sq_le}
    \leanok
    \uses{thm:hastings05_projection_deficiency,
        thm:hastings05_nested_projections,thm:hastings05_normalization,
        thm:hastings05_compressed_window_projection,
        thm:hastings05_compressed_window_bound}
    Let $P,P'$ be commuting orthogonal projections, $M=M^\dagger$, $c\in\R$,
    $x>0$, and $\psi$ a unit vector with $\lVert(M-c)\psi\rVert\le\epsilon$,
    $\lVert\psi-P\psi\rVert\le D$, and $\lVert\psi-P'\psi\rVert\le D$. Let $R$
    be the spectral window of $M$ on $[c-x,c+x]$ compressed to
    $\operatorname{ran}(PP')$ (Definition~\ref{def:hastings05_compressed_window}).
    Then
    \begin{align}
        \lVert\psi-R\psi\rVert^2
         & \le2D^2+\Bigl(\frac{\epsilon+2\lVert M-c\rVert D}{x}\Bigr)^2 .
        \label{eq:hastings05_level_bound}
    \end{align}
\end{lemma}
\begin{proof}\leanok
    Put $Q=PP'$ and $\psi'=Q\psi/\lVert Q\psi\rVert\in\operatorname{ran}Q$.
    By Lemma~\ref{thm:hastings05_projection_deficiency} and
    Lemma~\ref{thm:hastings05_normalization},
    $\lVert\psi-Q\psi\rVert^2\le2D^2$ and $\lVert\psi'-\psi\rVert\le2D$, so
    Lemma~\ref{thm:hastings05_compressed_window_bound} gives
    $\lVert\psi'-R\psi'\rVert\le(\epsilon+2\lVert M-c\rVert D)/x$. Since
    $RQ=R$, Lemma~\ref{thm:hastings05_nested_projections} yields
    (\ref{eq:hastings05_level_bound}).
\end{proof}

\begin{definition}[The deficiency bound of a level]
    \label{def:hastings05_hierarchy_deficiency}
    \lean{TNLean.MPS.AreaLaw.LocalSolution.hierarchyDeficiency,
        TNLean.MPS.AreaLaw.LocalSolution.hierarchyDeficiency_zero,
        TNLean.MPS.AreaLaw.LocalSolution.hierarchyDeficiency_succ}
    \leanok
    For $\epsilon,K\in\R$ and radii $x_0,x_1,\dots$, set $D_0=0$ and
    \begin{align}
        D_{m+1} & =\Bigl(2D_m^2+\Bigl(\frac{\epsilon+2KD_m}{x_m}\Bigr)^2
        \Bigr)^{1/2}.
        \label{eq:hastings05_hierarchy_deficiency}
    \end{align}
\end{definition}

\begin{lemma}[All levels]
    \label{thm:hastings05_hierarchy_levels}
    \lean{TNLean.MPS.AreaLaw.LocalSolution.isStarProjection_hierarchy,
        TNLean.MPS.AreaLaw.LocalSolution.norm_sub_hierarchy_le}
    \leanok
    \uses{thm:hastings05_level_against_ground_state,
        def:hastings05_hierarchy_deficiency}
    Let $P_k^0=1$, and let $P_k^{m+1}$ be the window of $M_k^m=(M_k^m)^\dagger$
    at $c_k^m$ with radius $x_m>0$ inside $\operatorname{ran}(P_{2k}^mP_{2k+1}^m)$.
    If $[P_{2k}^m,P_{2k+1}^m]=0$, $\lVert(M_k^m-c_k^m)\psi\rVert\le\epsilon$,
    and $\lVert M_k^m-c_k^m\rVert\le K$ for all $m,k$ and a unit vector $\psi$,
    then every $P_k^m$ is an orthogonal projection and
    $\lVert\psi-P_k^m\psi\rVert\le D_m$ (with $\psi=\Psi_a$ and
    $c_k^m=\langle\Psi_a,M_k^m\Psi_a\rangle$ in
    \cite[lines 668--707]{Hastings2006Solving}).
\end{lemma}
\begin{proof}\leanok
    Induct on $m$, applying Lemma~\ref{thm:hastings05_level_against_ground_state}
    to $P_{2k}^m,P_{2k+1}^m$ and comparing with
    (\ref{eq:hastings05_hierarchy_deficiency}) via
    $\lVert M_k^m-c_k^m\rVert\le K$.
\end{proof}

\begin{definition}[The majorant for the source's radii]
    \label{def:hastings05_hierarchy_majorant}
    \lean{TNLean.MPS.AreaLaw.LocalSolution.hierarchyMajorant,
        TNLean.MPS.AreaLaw.LocalSolution.hierarchyMajorant_zero,
        TNLean.MPS.AreaLaw.LocalSolution.hierarchyMajorant_succ}
    \leanok
    For $e,s\in\R$, set $\beta_0=0$ and
    $\beta_{m+1}=\bigl(2\beta_m^2+(e+2\beta_m)^{2s}\bigr)^{1/2}$.
\end{definition}

\begin{lemma}[The source's radii]
    \label{thm:hastings05_hierarchy_source_radii}
    \lean{TNLean.MPS.AreaLaw.LocalSolution.hierarchyMajorant_nonneg,
        TNLean.MPS.AreaLaw.LocalSolution.hierarchyDeficiency_le_hierarchyMajorant}
    \leanok
    \uses{def:hastings05_hierarchy_deficiency,def:hastings05_hierarchy_majorant}
    For $K,e>0$, $\epsilon=Ke$, and $x_m=K(e+2\beta_m)^{1-s}$
    \cite[lines 676--681]{Hastings2006Solving}, $0\le D_m\le\beta_m$ for all $m$.
\end{lemma}
\begin{proof}\leanok
    By induction on $m$ using~(\ref{eq:hastings05_hierarchy_deficiency}) and
    $(\epsilon+2KD_m)/x_m\le(e+2\beta_m)/(e+2\beta_m)^{1-s}=(e+2\beta_m)^s$.
\end{proof}

\begin{lemma}[The closed form]
    \label{thm:hastings05_hierarchy_closed_form}
    \lean{TNLean.MPS.AreaLaw.LocalSolution.hierarchyMajorant_le,
        TNLean.MPS.AreaLaw.LocalSolution.hierarchyMajorant_le_rpow_quarter}
    \leanok
    \uses{def:hastings05_hierarchy_majorant,thm:hastings05_exponent}
    For $0<e\le1$ and $0\le s\le1$, $\beta_m\le(3^m-1)e^{s^m}$. In particular,
    for $s=1-1/h$ and $h\ge2$ ($N=2^h$ supersites),
    $\beta_h\le(3^h-1)e^{1/4}$, in place of the printed
    $N\exp[-(l_{\mathrm{proj}}/l_{\mathrm{fac}})(1-1/h)^h]$ at
    $e=\exp(-l_{\mathrm{proj}}/l_{\mathrm{fac}})$
    \cite[lines 705--707]{Hastings2006Solving}.
\end{lemma}
\begin{proof}\leanok
    Set $C_m=3^m-1$. By induction on $m$, since $e\le e^{s^m}\le e^{s^{m+1}}$
    and $y^s\le y$ for $y\ge1$,
    \begin{align}
        (e+2\beta_m)^s\le(1+2C_m)e^{s^{m+1}},\qquad
        2C_m^2+(1+2C_m)^2\le(3C_m+2)^2=C_{m+1}^2.
        \notag
    \end{align}
    For $s=1-1/h$ and $m=h$, Lemma~\ref{thm:hastings05_exponent} gives
    $(1-1/h)^h\ge1/4$, so $e^{(1-1/h)^h}\le e^{1/4}$.
\end{proof}

\subsection{The supersite grouping}

Grouping sites into supersites of length $2l_{\mathrm{proj}}$ makes each $M_i$
act on at most two adjacent supersites, with $M_{k,k+1}$ the sum of $M_i$ on
supersites $k,k+1$~\cite[lines 609--628]{Hastings2006Solving}.

\begin{lemma}[Spectral windows keep supports]
    \label{thm:hastings05_window_support}
    \lean{TNLean.MPS.AreaLaw.LocalSolution.cfc_mem_supportedOperators,
        TNLean.MPS.AreaLaw.LocalSolution.windowProjection_mem_supportedOperators,
        TNLean.MPS.AreaLaw.LocalSolution.compressedWindowProjection_mem_supportedOperators}
    \leanok
    \uses{def:hastings05_window_projection,def:hastings05_compressed_window}
    If $A$ is supported on $S$, then $f(A)$ is supported on $S$ for every
    $f:\R\to\R$; in particular, the spectral window of $M$ and its compression
    to $\operatorname{ran}Q$ are supported on $S$ whenever $M$ and $Q$ are.
\end{lemma}
\begin{proof}\leanok
    Operators supported on $S$ form a closed $*$-subalgebra, which is preserved
    by the continuous functional calculus and products.
\end{proof}

\begin{definition}[Blocks of a level]
    \label{def:hastings05_level_block}
    \lean{TNLean.MPS.AreaLaw.LocalSolution.levelBlock}
    \leanok
    For a supersite map $\sigma:\{0,\dots,L-1\}\to\N$, block $k$ of level $m$
    is $B_k^m=\{i:\lfloor\sigma(i)/2^m\rfloor=k\}$.
\end{definition}

\begin{lemma}[Halves of a block]
    \label{thm:hastings05_level_block_halves}
    \lean{TNLean.MPS.AreaLaw.LocalSolution.levelBlock_two_mul_subset,
        TNLean.MPS.AreaLaw.LocalSolution.levelBlock_two_mul_add_one_subset,
        TNLean.MPS.AreaLaw.LocalSolution.disjoint_levelBlock_two_mul}
    \leanok
    \uses{def:hastings05_level_block}
    For all $m,k$, $B_{2k}^m\cap B_{2k+1}^m=\emptyset$ and
    $B_{2k}^m\cup B_{2k+1}^m\subseteq B_k^{m+1}$.
\end{lemma}
\begin{proof}\leanok
    Use $\lfloor\sigma(i)/2^{m+1}\rfloor=\lfloor\lfloor\sigma(i)/2^m\rfloor/2\rfloor$.
\end{proof}

\begin{lemma}[All levels on supersites]
    \label{thm:hastings05_supersite_hierarchy}
    \lean{TNLean.MPS.AreaLaw.LocalSolution.mem_supportedOperators_hierarchy,
        TNLean.MPS.AreaLaw.LocalSolution.commute_supersiteHierarchy,
        TNLean.MPS.AreaLaw.LocalSolution.norm_sub_supersiteHierarchy_le}
    \leanok
    \uses{thm:hastings05_hierarchy_levels,thm:hastings05_window_support,
        thm:hastings05_level_block_halves}
    In Lemma~\ref{thm:hastings05_hierarchy_levels}, if each $M_k^m$ is
    supported on $B_k^{m+1}$, then $P_k^m$ is supported on $B_k^m$,
    $[P_{2k}^m,P_{2k+1}^m]=0$, and $\lVert\psi-P_k^m\psi\rVert\le D_m$.
\end{lemma}
\begin{proof}\leanok
    By induction on $m$ using Lemmas~\ref{thm:hastings05_window_support}
    and~\ref{thm:hastings05_level_block_halves}, $P_{2k}^m$ and $P_{2k+1}^m$
    have disjoint supports $B_{2k}^m$ and $B_{2k+1}^m$, hence commute, and
    Lemma~\ref{thm:hastings05_hierarchy_levels} applies.
\end{proof}

\begin{lemma}[The residual of a sum]
    \label{thm:hastings05_sum_residual}
    \lean{TNLean.MPS.AreaLaw.LocalSolution.norm_toEuclideanLin_sum_sub_le}
    \leanok
    For operators $M_i$ and scalars $c_i$,
    $\lVert(\sum_iM_i-\sum_ic_i)\psi\rVert\le\sum_i\lVert(M_i-c_i)\psi\rVert$;
    for $c_i=\langle\Psi_a,M_i\Psi_a\rangle$ this bounds the residual of
    $M_{k,k+1}$ from eq.~(low)~\cite[lines 628--639]{Hastings2006Solving}.
\end{lemma}
\begin{proof}\leanok
    By the triangle inequality.
\end{proof}
